\subsection{Constructing a Common Unit Length}

\begin{derivationbox}[Choosing and transferring the unit]
Choose a positive segment length and call it one unit. Using a compass, transfer
this same length from $O$ repeatedly along both axes and in both directions. The
first marks are labelled $1$ and $-1$; repetition constructs the integer marks
\[
\ldots,-3,-2,-1,0,1,2,3,\ldots
\]
on each axis.
\end{derivationbox}

\begin{center}
\begin{tikzpicture}[scale=0.72]
    \draw[axis] (-4.5,0) -- (4.7,0) node[right] {$x$};
    \draw[axis] (0,-4.4) -- (0,4.6) node[above] {$y$};

    \draw[dashed,line width=0.8pt,draw=blue!70] (0,0) circle (1);

    \foreach \x/\col in {1/orange!85!black,-1/orange!85!black,
                           2/green!55!black,-2/green!55!black,
                           3/violet!75,-3/violet!75} {
        \draw[dashed,line width=0.8pt,draw=\col] (\x,0) circle (1);
    }
    \foreach \y/\col in {1/orange!85!black,-1/orange!85!black,
                           2/green!55!black,-2/green!55!black,
                           3/violet!75,-3/violet!75} {
        \draw[dashed,line width=0.8pt,draw=\col] (0,\y) circle (1);
    }

    \foreach \x in {-3,-2,-1,1,2,3} {
        \draw (\x,0.09) -- (\x,-0.09) node[below=3pt] {$\x$};
    }
    \foreach \y in {-3,-2,-1,1,2,3} {
        \draw (0.09,\y) -- (-0.09,\y) node[left=3pt] {$\y$};
    }
    \node[below left] at (0,0) {$0$};

    \fill[blue!70] (0,0) circle (1.8pt);
    \foreach \x in {1,-1} {\fill[orange!85!black] (\x,0) circle (1.8pt);}
    \foreach \y in {1,-1} {\fill[orange!85!black] (0,\y) circle (1.8pt);}
    \foreach \x in {2,-2} {\fill[green!55!black] (\x,0) circle (1.8pt);}
    \foreach \y in {2,-2} {\fill[green!55!black] (0,\y) circle (1.8pt);}
    \foreach \x in {3,-3} {\fill[violet!75] (\x,0) circle (1.8pt);}
    \foreach \y in {3,-3} {\fill[violet!75] (0,\y) circle (1.8pt);}

    \draw[<->,gray] (0,-4.0) -- (1,-4.0);
    \node[gray,below] at (0.5,-4.0) {one common unit};
\end{tikzpicture}
\end{center}

\begin{corebox}[Why the unit must be common]
A segment labelled as one unit on the $x$-axis must be congruent to a segment
labelled as one unit on the $y$-axis. Otherwise the two coordinate numbers would
measure lengths in incompatible scales.
\end{corebox}

\begin{interpretationbox}[Connection with distance]
The common scale is what later allows horizontal and vertical coordinate
differences to be combined by the Pythagorean theorem in the standard Euclidean
distance formula.
\end{interpretationbox}
